\subsection{外希尔伯特幂}

设 E 为希尔伯特空间，n 为正整数。对每个置换 \(\sigma\in\mathfrak{S}_{n}\)，以 \(\varepsilon_{\sigma}\) 表示其符号差；令 \(a_{n}=\frac{1}{n!}\sum_{\sigma\in\mathfrak{S}_{n}}\varepsilon_{\sigma}p_{\sigma}\)（在 \(\mathcal{L}(\hat{\mathbf{T}}^{n}(\mathrm{E}))\) 中）（V, p.~29）。立即可见 \(a_{n}\) 是一个正交投影算子，其像 \(\overline{\mathbf{A}_{n}^{\prime}(\mathrm{E})}\) 是 \(\hat{\mathbf{T}}^{n}(\mathrm{E})\) 中空间 \(\mathbf{A}_{n}^{\prime}(\mathrm{E})\)（由所有 \(n(\mathrm{A},\mathrm{III},\S7,\mathrm{No}.4)\) 阶反对称张量组成的）的闭包。存在一个同构 \(\pi_{n}\)（从 \(\Lambda^{n}(\mathrm{E})\) 到 \(\mathbf{A}_{n}^{\prime}(\mathrm{E})\) 上），它由下式刻画

\[
\pi_ {n} (x _ {1} \wedge \ldots \wedge x _ {n}) = \mathbf {a} _ {n} (x _ {1} \otimes \ldots \otimes x _ {n}) = \frac {1}{n !} \sum_ {\sigma \in \mathfrak {S} _ {n}} \varepsilon_ {\sigma} x _ {\sigma (1)} \otimes \ldots \otimes x _ {\sigma (n)}\tag{24}
\]

其中 \(x_{1}, \ldots, x_{n}\) 属于 E。现在我们可以通过令下式成立而在 \(\Lambda^{n}(E)\) 上定义一个分离准希尔伯特空间结构

\[
\langle u | v \rangle = n! \left\langle \pi_ {n} (u) \mid \pi_ {n} (v) \right\rangle .\tag{25}
\]

更明确地，我们有（试与 A, III, § 11, No.~5 的公式 (30) 比较）。

\[
\langle x _ {1} \wedge \dots \wedge x _ {n} | y _ {1} \wedge \dots \wedge y _ {n} \rangle = \det (\langle x _ {i} | y _ {j} ^ {\prime} \rangle)\tag{26}
\]

其中 \(x_{1},\ldots ,x_{n}\) 与 \(y_{1},\ldots ,y_{n}\) 都属于 E。

设 \(\hat{\mathbf{A}}^n (\mathrm{E})\) 表示准希尔伯特空间 \(\mathbf{A}^n (\mathrm{E})\) 的完备化，\(\hat{\mathbf{A}} (\mathrm{E})\) 表示诸希尔伯特空间 \(\hat{\mathbf{A}}^n (\mathrm{E})\) 的外希尔伯特和。

例. --- * 采用 V, p.~32 例 1 的记号，我们可以把希尔伯特空间 \(\symbf{\Lambda}(\mathrm{E})\) 与所有可测函数序列 \((f_n)_{n\geqslant 0}\)（对它们，(22) 中定义的数 \(\| \mathbf{f}\|\) 为有限）组成的集合等同起来，其中每个函数 \(f_{n}\) 都是反对称的，即满足关系式

\[
f _ {n} (\mathbf {x} _ {\sigma (1)}, \dots , \mathbf {x} _ {\sigma (n)}) = \varepsilon_ {\sigma} f _ {n} (\mathbf {x} _ {1}, \dots , \mathbf {x} _ {n})
\]

其中 \(\sigma \in \mathfrak{S}_n\) 为任意置换。希尔伯特空间 \(\hat{\Lambda}(\mathrm{E})\) 称为对应于权 \(\omega\) 的反对称 Fock 空间。*

命题 5. --- 设 \((e_i)_{i \in I}\) 为希尔伯特空间 E 的标准正交基。给 I 赋以一个全序结构。于是所有元素 \(e_{i_1} \wedge \ldots \wedge e_{i_n}\)（对 \(i_1 < \cdots < i_n\)）组成的集合是 \(\hat{\Lambda}^n(E)\) 的标准正交基。

我们知道（A, III, § 7, No.~8），所说的这些元素构成向量空间 \(\Lambda^n (\mathrm{E}_0)\) 的基，其中 \(\mathrm{E}_0\) 是 E 中由向量 \(e_i\) 生成的向量子空间。再者，当 \(i_1 < \dots < i_n\) 时，内积 \(\langle e_{i_k}|e_{i_\ell}\rangle\) 组成的矩阵是 \(n\) 阶单位矩阵；由 (26)，有 \(\| e_{i_1}\wedge \ldots \wedge e_{i_n}\| = 1\)。最后，若 \((i_1,\dots,i_n)\) 与 \((j_1,\dots,j_n)\) 是 I 中元素的两个不同的严格递增序列，则存在一个元素 \(j_\ell\) 与 \(i_1,\ldots ,i_n\) 都不同，于是 \(\langle e_{i_k}|e_{j_\ell}\rangle = 0\)（对 \(1\leqslant k\leqslant n\)），从而由 (26) 得 \(\langle e_{i_1}\wedge \ldots \wedge e_{i_n}|e_{j_1}\wedge \ldots \wedge e_{j_n}\rangle = 0\)。换言之，元素 \(e_{i_1}\wedge \ldots \wedge e_{i_n}\)（对 \(i_1 < \dots < i_n\)）组成的族是标准正交的。

但 \(E_{0}\) 在 E 中稠密，且映射 \((x_{1},...,x_{n})\mapsto x_{1}\wedge\ldots\wedge x_{n}\)（从 \(E\times\cdots\times E\) 到 \(\Lambda^{n}(E)\) 中）是连续的。因此 \(\Lambda^{n}(E_{0})\) 在 \(\Lambda^{n}(E)\) 中稠密，由此即得命题 5。

推论. --- 假定希尔伯特空间 E 是两个正交子空间 M 与 N 的直和。从 \(\symbf{\Lambda}(\mathbf{M})\otimes\symbf{\Lambda}(\mathbf{N})\) 到 \(\symbf{\Lambda}(\mathbf{E})\) 上的典范同构 g（A, III, § 7, No.~7）以唯一的方式延拓为从 \(\hat{\symbf{\Lambda}}(\mathbf{M})\hat{\otimes}_{2}\hat{\symbf{\Lambda}}(\mathbf{N})\) 到 \(\hat{\symbf{\Lambda}}(\mathbf{E})\) 上的希尔伯特空间同构。

证明与命题 4 的推论（V, p.~31）的证明类似。

设 \(\mathbf{E}\) 与 \(\mathbf{F}\) 为两个希尔伯特空间，\(u \in \mathcal{L}(\mathbf{E};\mathbf{F})\)。与对称幂 \(\hat{\mathbf{S}}^n (\mathbf{E})\) 的情形（V, p.~32）一样，我们将证明线性映射 \(\symbf{\Lambda}^n (u)\)（从 \(\symbf{\Lambda}^n (\mathbf{E})\) 到 \(\symbf{\Lambda}^n (\mathbf{F})(\mathbf{A},\mathbf{III},\S 7,\mathbf{No}.4)\) 中）可延拓为连续线性映射 \(\hat{\symbf{\Lambda}}^n (u)\)（从 \(\hat{\symbf{\Lambda}}^n (\mathbf{E})\) 到 \(\hat{\symbf{\Lambda}}^n (\mathbf{F})\) 中）。我们有关系式

\[
\hat {\symbf {\Lambda}} ^ {n} (1 _ {\mathrm{E}}) = 1 _ {\hat {\symbf {\Lambda}} ^ {n} (\mathrm{E})},\tag{27}
\]

\[
\hat {\symbf {\Lambda}} ^ {n} (v \circ u) = \hat {\symbf {\Lambda}} ^ {n} (v) \circ \hat {\symbf {\Lambda}} ^ {n} (u) \quad \text { if } v \text { belongs   to } \mathscr {L} (\mathrm{F}; \mathrm{G}),\tag{28}
\]

\[
\left\| \hat {\symbf {\Lambda}} ^ {n} (u) \right\| \leqslant \| u \| ^ {n}.\tag{29}
\]

一般而言，公式 (29) 中的等号并不成立（TS, IV, § 6）。最后，我们有一个同构 \(\psi_{n} = \psi_{n,\mathrm{E}}\)（从 \(\hat{\mathbf{A}}^{n}(\mathrm{E})\) 到子空间 \(\overline{\mathbf{A}_{n}^{\prime}(\mathrm{E})}\)（\(\hat{\mathbf{T}}^{n}(\mathrm{E})\) 的）上），它由下式定义

\[
\psi_ {n} (x _ {1} \wedge \dots \wedge x _ {n}) = \frac {1}{(n !) ^ {1 / 2}} \sum_ {\sigma \in \mathfrak {S} _ {n}} \varepsilon_ {\sigma} x _ {\sigma (1)} \otimes \dots \otimes x _ {\sigma (n)}.\tag{30}
\]

